\subsection{Proof of Theorem \ref{thm:metrics}}\label{proof:theorem1}
Suppose that \eqref{eq:rprs} holds. This implies,
\begin{align*}
	 & \ip{p_T \Psi\bmat{G_\Delta \\ I}w_\Delta}{p_T\mcl{M}_\mcl{V}\Psi\bmat{G_\Delta \\ I}w_\Delta}_{\mbf L_{[a, b]}} \\ &+ \ip{p_T \bmat{G_{z\Delta} \\ 0}w_\Delta}{p_T\mcl{M}_\mcl{\hat{V}}\bmat{G_{z\Delta} \\ 0}w_\Delta}_{\mbf L_{[a, b]}}\\ &\leq -\epsilon \|p_Tw_\Delta\|^2_{\mbf L_{[a, b]}}
\end{align*}
Since $\mcl{\hat{V}}_{11}\succeq0$, we can then infer,
\begin{align*}
	 & \ip{p_T \Psi\bmat{G_\Delta  w_\Delta\\ w_\Delta}}{p_T\mcl{M}_\mcl{V}\Psi\bmat{G_\Delta w_\Delta \\ w_\Delta}}_{\mbf L_{[a, b]}}\!\!\!\!\!\!\!\!\!\!\!\!\leq -\epsilon \|p_Tw_\Delta\|^2_{\mbf L_{[a, b]}}
\end{align*}
Then, using Theorem 6 in \cite{10156335}, we find that \textbf{robust stability} is guaranteed with respect to performance multiplier $\mcl{M}_{\hat{V}}$. Moreover, rewriting \eqref{eq:rprs} as,
\begin{align*}
	 & \ip{p_T\Psi\bmat{z_\Delta\\w_\Delta}}{p_T\mcl{M}_{{\mcl{V}}}\Psi\bmat{z_\Delta \\ w_\Delta}}_{\mbf L_{[a, b]}} + \\&\ip{p_T\bmat{z\\w}}{p_T\mcl{M}_{\hat{\mcl{V}}}\bmat{z \\ w}}_{\mbf L_{[a, b]}}
	\!\!\!\!\!\!\!\!\!\leq  -\epsilon (\|p_T{w_\Delta}\|^2_{\mbf L_{[a, b]}}+\|p_T w\|^2_{\mbf L_{[a, b]}}),
\end{align*}
and given that $\Psi$ and $\mcl{M}_\mcl{V}$ satisfy the IQC in Definition \ref{def:multipliers} leads to,
\begin{align*}
	\ip{p_T\bmat{z\\w}}{p_T\mcl{M}_{\hat{\mcl{V}}}\bmat{z \\ w}}_{\mbf L_{[a, b]}}\!\!\!\!\!\!\!\!\!\leq  -\epsilon \|p_T w\|^2_{L_2}.
\end{align*}
Thus by Definition \ref{def:robperf} \textbf{robust performance} is achieved.
\subsection{Proof of Theorem \ref{thm:dualscaledIO}}\label{proof:theorem2}

\begin{sublemma}\label{lem:scalingisPI}
	The \textbf{Primal I/O Scaling} is a PI operator and invertible if and only if the \textbf{Dual I/O Scaling} is a PI operator and invertible
\end{sublemma}
\begin{proof}
	To show sufficiency, let $\bar{\mcl{S}} = (\hat{\mcl{J}}^*\mcl{S}\hat{\mcl{J}})^{-*} = \hat{\mcl{J}}^*\mcl{S}^{-*}\hat{\mcl{J}} = \hat{\mcl{J}}^*(\mcl{S}^{-1})^{*}\hat{\mcl{J}}$, where $\hat{\mcl{J}} = \bmat{0 & -I \\ I & 0} \in \Pi_4$. Then if $\mcl{S}$ is invertible then $\mcl{S}^{-1} \in \Pi_4$ and by using that $\Pi_4$ forms a composition algebra we find that the \textbf{Dual I/O Scaling} is PI and invertible. For the necessary case, notice $\mcl{S} = (\hat{\mcl{J}}^*\bar{\mcl{S}}\hat{\mcl{J}})^{-*}$.
\end{proof}
\begin{sublemma}
	If $\mcl{S}_{11} - \mcl{D}\mcl{S}_{21}$ and $\bar{\mcl{S}}_{11} - \mcl{D}^*\bar{\mcl{S}}_{21}$ are invertible, then the \textbf{Primal Scaled System} and \textbf{Dual Scaled System} are well-posed, respectively.
\end{sublemma}
\begin{proof}
	Applying,
	\begin{equation*}
		\bmat{z(t)\\w(t)} = \bmat{\mcl{S}_{11} & \mcl{S}_{12} \\ \mcl{S}_{21} & \mcl{S}_{22}}\bmat{\tilde{z}(t) \\ \tilde{w}(t)},
	\end{equation*}
	on,
	\[
	\begin{aligned}
			\partial_t(\mcl{T}\mbf{x})(t)                         & = \mcl{A}\mbf{x}(t) + \mcl{B}{w}(t), \\
			{z}(t) & = \mcl{C} \mbf{x}(t) + \mcl{D}{w}(t).
		\end{aligned}
	\]
	 yields,
	\[
		\begin{aligned}
			\partial_t(\mcl{T}\mbf{x})(t)                         & = \mcl{A}\mbf{x}(t) + \mcl{B}({\mcl{S}}_{21}\tilde{z}(t) + {\mcl{S}}_{22}\tilde{w}(t)), \\
			({\mcl{S}}_{11} - \mcl{D}{\mcl{S}}_{21}) \tilde{z}(t) & = \mcl{C} \mbf{x}(t) + (\mcl{D}{\mcl{S}}_{22} - {\mcl{S}}_{12}) \tilde{w}(t).
		\end{aligned}
	\]
	Then the \textbf{Primal Scaled System} is well-posed if and only if \( \mcl{S}_{11} - \mcl{D}\mcl{S}_{21} \) is invertible, ensuring a unique solution for \( \tilde{z}(t) \) in terms of \( \mbf{x}(t) \) and \( \tilde{w}(t) \). A similar proof shows that the \textbf{Dual Scaled System} is well-posed if $\bar{\mcl{S}}_{11} - \mcl{D}^*\bar{\mcl{S}}_{21}$ is invertible.
\end{proof}
% \begin{sublemma}
% 	Under the assumptions above, the dynamically scaled primal and dual systems are well-posed.
% \end{sublemma}
% \begin{proof}
% 	Given the dynamic multipliers,
% 	\begin{equation*}
% 		\bmat{\tilde{z}_\Delta \\ \tilde{w}_\Delta} = \bmat{\psi_{11} & \psi_{12} \\ \psi_{21} & \psi_{22}}\bmat{{z}_\Delta(t) \\ {w}_\Delta(t)},
% 	\end{equation*}
% 	which can be equivalently expressed as
% 	\begin{align*}
% 		\mcl{T}_\Psi \dot{x}_\Psi(t) & = \mcl{A}_\Psi x_\Psi(t) + \mcl{B}_\Psi \bmat{{z}_\Delta(t) \\ {w}_\Delta(t)},\\
% 		\bmat{\tilde{z}_\Delta(t) \\ \tilde{w}_\Delta(t)} &= \mcl{C}_\Psi x_\Psi(t) + \mcl{D}_\Psi\bmat{{z}_\Delta(t) \\ {w}_\Delta(t)}.
% 	\end{align*}
% 	Under the given conditions we can find the system,
% 	\begin{align*}
% 		\mcl{T}_? \dot{x}_?(t) &= \mcl{A}_? x_?(t) + \mcl{B}_? \tilde{w}_\Delta(t),
% 		\\ \tilde{z}_\Delta(t) &= \mcl{C}_? x_?(t) + \mcl{D}_? \tilde{w}_\Delta(t),
% 	\end{align*}
% \end{proof}
To show sufficiency, let $\mbf{x}(t)$, $w(t)$, and $z(t)$ satisfy the primal system dynamics with $\mbf{x}(0)=0$, and let $\tilde{z}, \tilde{w}$ be defined by the primal scaling in 1):
\begin{equation}\label{eqn:dual_scaling_explicit}
	{z} = \mcl{S}_{11}\tilde{z} + \mcl{S}_{12}\tilde{w},\quad w = \mcl{S}_{21}\tilde{z} + \mcl{S}_{22}\tilde{w}.
\end{equation}
From the proof of Theorem 9 in \cite{shivaknew}, specifically (12), we established that the unscaled signals, $z(t)$, $w(t)$, $\bar{z}(t)$ and $\bar{w}(t)$, satisfy the identity:
\begin{equation}\label{eq:identity}
	\int_0^t\ip{\bar{z}(\theta)}{w(t-\theta)}d\theta = \int_0^t\ip{\bar{w}(\theta)}{z(t-\theta)}d\theta.
\end{equation}
Substituting the dual scaling expressions from Eq.~\eqref{eqn:dual_scaling_explicit} into the left and right sides respectively, we obtain
\begin{equation*}
	\begin{aligned}
		 & \int_0^t\ip{\bar{z}(\theta)}{\mcl{S}_{21}\tilde{z}(t-\theta) + \mcl{S}_{22}\tilde{w}(t-\theta)}d\theta    \\
		 & = \int_0^t\ip{\bar{w}(\theta)}{\mcl{S}_{11}\tilde{z}(t-\theta) + \mcl{S}_{12}\tilde{w}(t-\theta)}d\theta.
	\end{aligned}
\end{equation*}
Rearranging the terms to group by $\tilde{\bar{z}}$ and $\tilde{\bar{w}}$ yields
\begin{equation*}
	\begin{aligned}
		 & \int_0^t \ip{\mcl{S}^*_{22}\bar{z}(\theta) - \mcl{S}^*_{12}\bar{w}(\theta)}{\tilde{w}(t-\theta)} d\theta     \\
		 & = \int_0^t \ip{-\mcl{S}^*_{21}\bar{z}(\theta) + \mcl{S}^*_{11}\bar{w}(\theta)}{\tilde{z}(t-\theta)} d\theta.
	\end{aligned}
\end{equation*}
We observe that the terms inside the inner product match the definition of the dual scaled variables:
\begin{equation*}
	\tilde{\bar{w}} = -\mcl{S}_{21}^*\bar{z} + \mcl{S}_{11}^*\bar{w}, \quad \tilde{\bar{z}} = \mcl{S}_{22}^*\bar{z} - \mcl{S}^*_{12}\bar{w}.
\end{equation*}
Substituting these identities, we arrive at the scaled identity:
\begin{equation}\label{eqn:scaled_equivalence}
	\int_0^t\ip{\tilde{\bar{z}}(\theta)}{\tilde{w}(t-\theta)} d\theta = \int_0^t\ip{\tilde{\bar{w}}(\theta)}{\tilde{z}(t-\theta)} d\theta.
\end{equation}
This equation is structurally identical to \eqref{eq:identity}, but relates the scaled primal and dual variables. Now, we proceed with the truncation argument. For any fixed $\tau > 0$, let $\tilde{\bar{w}}(t)\in\RL$ be arbitrary. Under the assumption that the primal and dual scaled system are well-posed, i.e. {$\mcl{S}_{11} - \mcl{D}\mcl{S}_{21}$ and $\bar{\mcl{S}}_{11} - \mcl{D}^*\bar{\mcl{S}}_{21}$} are invertible, the scaled primal input can be defined as $\tilde{w}(t) = \tilde{\bar{z}}(\tau - t)$ for $t \leq \tau$ and $\tilde{w}(t) = 0$ for $t > \tau$. Then
\begin{equation*}
	\begin{aligned}
		\| P_{\tau}\tilde{\bar{z}}(t) \|_{\RL}^2 & = \int_{0}^{\tau} \ip{\tilde{\bar{z}}(s)}{\tilde{w}(\tau-s)} ds              \\&= \int_{0}^{\tau} \ip{\tilde{\bar{w}}(s)}{\tilde{z}(\tau-s)} ds \\
		                                         & \leq \| P_{\tau}\tilde{\bar{w}}(t) \|_{\RL} \| P_{\tau}\tilde{z}(t) \|_{\RL} \\&\leq \| P_{\tau}\tilde{\bar{w}}(t) \|_{\RL} \| P_{\tau}\tilde{w}(t) \|_{\RL}.
	\end{aligned}
\end{equation*}
Since by construction $\| P_{\tau}\tilde{w}(t) \|_{\RL} = \| P_{\tau}\tilde{\bar{z}}(t) \|_{\RL}$, thus we conclude
\begin{equation*}
	\| P_{\tau}\tilde{\bar{z}}(t) \|_{\RL} \leq  \| P_{\tau}\tilde{\bar{w}}(t) \|_{\RL}
\end{equation*}
Since $\mcl{T}^{**} = \mcl{T}$, Sublemma \ref{lem:scalingisPI}, and
\begin{equation*}
	\bmat{(\mcl{A}^*)^* & (\mcl{B}^*)^*\\(\mcl{C}^*)^*&(\mcl{D}^*)^*}=
	{\bmat{\mcl{A}^* & \mcl{C}^*\\\mcl{B}^*&\mcl{D}^*}}^*=
	\bmat{\mcl{A} & \mcl{B}\\\mcl{C}&\mcl{D}}^{**}=\bmat{\mcl{A} & \mcl{B}\\\mcl{C}&\mcl{D}}
\end{equation*}we have that sufficiency implies necessity.
\subsection{Proof of Lemma \ref{lem:PNPIfact}}\label{proof:PNPIfact}
Given the operators $\mcl Q, \mcl{S}, \mcl{R} \in \Pi_4$
\begin{equation*}
	\mcl V=\bmat{\mcl Q & \mcl S \\ \mcl S^* & \mcl R}.
\end{equation*}
Since $\mcl V$ is a strict PN-PI operator, there exists $\epsilon>0$ such that
\begin{equation*}
	\mcl Q \ge \epsilon I, \qquad \mcl R\le -\epsilon I.
\end{equation*}
In particular, $\mcl Q$ is self-adjoint and coercive. Hence, by the positive
square-root property for coercive self-adjoint PI operators, there
exists a unique positive invertible operator $\mcl S_{11}$ such that
\begin{equation*}
	{\mcl Q=\mcl S_{11}^*\mcl S_{11}}, \quad
	\mcl S_{12}:=\mcl S_{11}^{-*}\mcl S, \quad
	\mcl S_{11}^*\mcl S_{12}=\mcl S.
\end{equation*}
Next define
\begin{equation*}
	\mcl W:=\mcl S_{12}^*\mcl S_{12}-\mcl R=\mcl S^*\mcl Q^{-1}\mcl S-\mcl R.
\end{equation*}
Since $\mcl S^*\mcl Q^{-1}\mcl S\ge 0$ and $-\mcl R\ge \epsilon I$, it follows that
\begin{equation*}
	\mcl W\ge \epsilon I.
\end{equation*}
Hence $\mcl W$ is self-adjoint, coercive, and positive. Again by the positive
square-root property, there exists a unique positive invertible operator
$\mcl S_{22}$ such that
\begin{equation*}
	\mcl W=\mcl S_{22}^*\mcl S_{22}.
\end{equation*}
Therefore,
\begin{equation*}
	\mcl R=\mcl S_{12}^*\mcl S_{12}-\mcl S_{22}^*\mcl S_{22}.
\end{equation*}
Combining these identities gives
\begin{equation*}
	\mcl V=\bmat{\mcl S_{11} & \mcl S_{12} \\ 0 & \mcl S_{22}}^*\bmat{I & 0 \\ 0 & -I}\bmat{\mcl S_{11} & \mcl S_{12} \\ 0 & \mcl S_{22}}.
\end{equation*}
By construction, $\mcl S_{11}>0$ and $\mcl S_{22}>0$.

\subsection{Proof of Lemma \ref{lem:primdualPNPI}}\label{proof:primdualPNPI}
Given the operators $\mcl Q, \mcl{S}, \mcl{R} \in \Pi_4$, write
\begin{equation*}
	\mcl K:=\bmat{-\mcl R & \mcl S^* \\ \mcl S & -\mcl{Q}},
	\qquad
	\bar{\mcl V}=\mcl K^{-1}.
\end{equation*}
Since \(\mcl V\) is a strict PN-PI operator, there exists \(\epsilon>0\) such that
\begin{equation*}
	\mcl Q\succeq \epsilon I,
	\qquad
	\mcl R\preceq -\epsilon I.
\end{equation*}
so both \(-\mcl R\) and \(\mcl Q\) are self-adjoint, coercive, and invertible. Applying the block inverse formula to \(\mcl K\), generalized to the PI-operator setting from \cite{lu2002inverses}, gives the following expressions. In addition, \cite[Theorem 2]{peet2018convexsolutionhinftyoptimalcontroller} shows that the inverse of a coercive, self-adjoint PI operator is again a coercive, self-adjoint PI operator.
\begin{equation*}
	\bar{\mcl Q}
	=
	\left(\mcl S^*\mcl Q^{-1}\mcl S -\mcl R\right)^{-1}
	{\succeq \epsilon I},
\end{equation*}
and
\begin{equation*}
	\bar{\mcl R}
	=
	-\left(\mcl Q+\mcl S(-\mcl R)^{-1}\mcl S^*\right)^{-1}
	{\preceq -\epsilon I}.
\end{equation*}
Moreover, by composition algebra $\bar{\mcl S}$ is also a PI operator,
\begin{equation*}
	\bar{\mcl S} = {\mcl Q}^{-1} \mcl{S}\bar{\mcl{Q}}
\end{equation*}
Hence \(\bar{\mcl V}\) is a strict PN-PI operator. The converse follows by the same argument with \(\mcl V\) and \(\bar{\mcl V}\) interchanged, since
\begin{equation*}
	\mcl V=\bmat{-\bar{\mcl R} & \bar{\mcl S}^* \\ \bar{\mcl S} & -\bar{\mcl Q}}^{-1}.
\end{equation*}
Thus \(\mcl V\) is a strict PN-PI operator if and only if \(\bar{\mcl V}\) is a strict PN-PI operator.

